% sec-00: front matter.  Abstract and the reading path.
\section*{Abstract}
\addcontentsline{toc}{section}{Abstract}

The National Science Foundation has asked whether ChemicalQDevice would apply to
a multi-million-dollar pilot program. In September 2026 the agency announced a
\$20 million, two-year Commercialization Readiness Pilot, run by the National
Commercialization and Translation Institute at the University of Central
Florida, that will select six to eight companies for milestone-based awards of
\$1 million to \$3.75 million each \cite{nsfncti2026,nctiucf2026}. Its competition
is open to active or recent NSF SBIR and STTR Phase II and IIB awardees. The
agency's Strategic Breakthrough tier, of up to \$30 million with a one-to-one
match, is likewise open only to Phase II awardees \cite{nsf26510}. The company
holds no NSF award, and so meets neither gate today.

This packet answers the inquiry exactly. It thanks the agency, states the gate
and the company's place against it, asks whether the inquiry contemplates a
route the company has missed, and submits the Project Pitch that is the first
step on the published route: Pitch, Phase I, Phase II, and only then the pilot.
NSF's SBIR program funds technical innovation and does not fund clinical
trials, so the Pitch proposes the part of the program NSF can fund: a governed,
advisory-only language model layer for robotic oncologic surgery, with
verification by independent models and a public audit trail. The robotic
Phase 1 itself stays with the National Institutes of Health and private capital,
and the packet shows how the program is split so that no dollar is asked for
twice.

\subsection*{How to read this}

A reader with five minutes takes \S1, the inquiry, the gate and the funding
ladder, and \S2, the day's actions with the path from Pitch to proposal.

A reviewer asking what NSF is being asked to fund takes \S3, the split between
NSF, NIH and private capital, and \S4, the four Phase I objectives, each with a
measure and a threshold.

A funder takes \S5, the merit review criteria against the program's evidence,
including the quantitative record behind the robotic Phase 1. \S6 carries the
positioning constraints, the method, and the references, every one of which is
a clickable link.
